% ARCHIVAL DRAFT: not included in the active manuscript. Current results are in the main source and generated result tables.
\section{Measured results and current scope}
\label{m:results-summary}

We evaluate a fixed F2 step-14,000 payload with 4,091,581,440 stored elements,
independently parameterized forward and backward mixers, and no depth-loop or latent
conditioner. Its digest and tensor inventory identify the evaluated artifact; its
correspondence to the cited public Hub revision remains unresolved. The payload stores
FP32 tensors and is evaluated in BF16. Detailed artifact qualification, complete tables,
and execution records appear in Sections~\ref{m:identity}--\ref{m:notmeasured}.

\paragraph{Capability under causal option scoring.}
Table~\ref{tab:main-capability} summarizes complete capability panels. Adapted checkpoints
use the causal \texttt{fwdce} scoring path and released checkpoints use \texttt{raw};
both compute the same length-normalized option-span likelihood statistic. MMLU uses
one-token answer labels. These evaluations measure a causal inference path through the
adapted weights, not iterative masked denoising or bidirectional generation. They
therefore characterize retained capability without establishing the proposed
long-context refinement benefit. Tokenizer-specific prompt trees fix the input
convention; RACE, OpenBookQA, and MMLU-Redux additionally pass exact unique-document
coverage checks. Upstream harness fields marked \texttt{unpinned} remain provenance
limitations despite the subsequently recorded source hashes.

\begin{table}[t]
\centering\small
\setlength{\tabcolsep}{4pt}
\begin{tabular}{@{}lrrrrr@{}}
\toprule
Model & MMLU & HellaSwag & RACE & OBQA & Redux\\
Items & 14,042 & 10,042 & 4,887 & 500 & 5,700\\
\midrule
F2 (4.09B) & 46.9 & 71.0 & 46.90 & 35.80 & 38.67\\
RWKV-7 2.9B & 54.6 & 74.3 & 47.64 & 41.20 & 42.74\\
RWKV-7 1.5B & 42.0 & 68.6 & 45.98 & 39.40 & 40.21\\
RWKV-7 0.4B & 26.0 & 53.4 & 40.58 & 34.80 & 33.95\\
\bottomrule
\end{tabular}
\caption{Complete-panel accuracy (\%); model sizes are total F2 stored elements or
released-model labels. OBQA is OpenBookQA and Redux is MMLU-Redux. MMLU/Redux use test
splits; the other panels use validation splits. Training histories and parameter counts
are unmatched. Full precision, uncertainty, and further comparators are in the detailed
results; no cross-dataset aggregate is formed.}
\label{tab:main-capability}
\end{table}

F2 trails released RWKV-7 2.9B on all five panels, including differences of
$-7.7$ percentage points on MMLU and $-3.3$ on HellaSwag. Paired-document bootstrap
intervals for F2 minus that release are $-0.74$ points $[-1.74,0.25]$ on RACE,
$-5.40$ $[-9.00,-2.00]$ on OpenBookQA, and $-4.07$ $[-5.04,-3.09]$ on
MMLU-Redux. These intervals condition on the evaluated checkpoints and executions;
they do not isolate architecture, objective, or pretraining exposure. F2 exceeds the
0.4B release on every panel, while its comparison with the 1.5B release is mixed.
MMLU-Redux is a separate benchmark, not a replacement for full MMLU.

\paragraph{Controlled direction and tied-depth adaptation.}
A separate matrix adapts a common 0.4B initialization on the same corpus and approximately
2.0B tokens, with three training seeds per cell (Table~\ref{tab:main-direction}).
Forward-only arms average 1.04 points higher on MMLU and 0.88 points higher on HellaSwag
than bidirectional arms under causal option scoring. Tied-loop mean differences are
at most 0.05 points. These observations establish neither a general direction penalty
nor a general loop null. Every adapted arm falls below its initialization on HellaSwag;
an equal-budget causal-continuation control would be needed to attribute that change
to the diffusion objective.

\begin{table}[t]
\centering\small
\begin{tabular}{@{}llrr@{}}
\toprule
Direction & Tied loop & MMLU & HellaSwag\\
\midrule
Released initialization & --- & 26.04 & 53.38\\
Forward-only & absent & 26.86 & 51.92\\
Forward-only & present & 26.89 & 51.95\\
Bidirectional & absent & 25.84 & 51.08\\
Bidirectional & present & 25.83 & 51.03\\
\bottomrule
\end{tabular}
\caption{Controlled 0.4B adaptation: accuracy (\%), means over three training seeds.
The initialization is a single released checkpoint. Seed ranges and interpretation
limits are reported in Section~\ref{m:direction}.}
\label{tab:main-direction}
\end{table}

\paragraph{Execution sensitivity, resources, and open endpoints.}
Repeated MMLU executions disagree on 0.20\% of F2 item verdicts and 0.75\% for two
looped checkpoints. Opposite changes can cancel in aggregate accuracy; these fractions
are not a universal resolution floor. Exploratory 4K--64K full-context forward probes
show approximately flat processed-input throughput for recurrent mixers and a roughly
$2\times$ decrease for transformers. The BiRWKV probe lacks resolved checkpoint
identity, so its memory figures are not F2 measurements. Neither input throughput nor
configured NFE establishes generation speed or deployment goodput.

The primary long-context study is underway. Its decoder must be selected on the
separate development panel and frozen before confirmation; no confirmation accuracy is
reported here. The theory specifies the law and approximation terms such an experiment
must distinguish. The current empirical results establish capability and measurement
boundaries, without claiming long-context superiority.
% Reserved integration point: add the frozen-decoder long-context table here only
% after complete, provenance-checked confirmation measurements are available.
